% Copy these examples into main.tex; the theme supplies the appearance.
\begin{frame}[fragile]{Start with an ordinary Beamer file}
  \begin{columns}[T,onlytextwidth]
    \begin{column}{.56\textwidth}
      \small
\begin{verbatim}
\documentclass[aspectratio=169]{beamer}
\usetheme{WVU}
\title{Presentation title}
\author{Your name}
\institute{West Virginia University}
\date{\today}

\begin{document}
\maketitle
% Your frames go here.
\thanksframe{Questions?}
\end{document}
\end{verbatim}
    \end{column}
    \begin{column}{.40\textwidth}
      \begin{enumerate}
        \item Open \texttt{main.tex}.
        \item Choose \textbf{LuaLaTeX}.
        \item Update your details.
        \item Replace the sample slides.
      \end{enumerate}
      \medskip
      Keep \texttt{beamerthemeWVU.sty} and the \texttt{styles/} folder
      beside your document.
      \medskip

      Fonts, bullets, boxes, and backgrounds are already set up.
    \end{column}
  \end{columns}
\end{frame}

\begin{frame}[fragile]{Keep Amurmaple's contact and layout features}
  \begin{columns}[T,onlytextwidth]
    \begin{column}{.50\textwidth}
      \textbf{Your details in the preamble}
      {\small
\begin{verbatim}
\mail{you@mail.wvu.edu}
\webpage{https://example.org}
\collaboration{Your collaboration}
\end{verbatim}
      }
      Email appears on the title and along the right edge of content
      slides. The webpage appears on the title; collaboration appears
      below the title and subtitle.
      \medskip

      Replace the starter's examples with your own details. Leave a field
      empty if you do not need it.
    \end{column}
    \begin{column}{.46\textwidth}
      \textbf{Familiar layout options}
      \begin{description}
        \item[\texttt{nogauge}] Hide the progress gauge.
        \item[\texttt{nomail}] Hide only the margin email.
        \item[\texttt{toplogo}] Put the sidebar logo at the top.
        \item[\texttt{leftframetitle}] Align slide titles left.
        \item[\texttt{sidebarwidth}] Set the sidebar width.
      \end{description}
      {\small
\begin{verbatim}
\usetheme[toplogo,
  sidebarwidth=64pt]{WVU}
\end{verbatim}
      }
      Replace your existing theme line when choosing options.
    \end{column}
  \end{columns}
\end{frame}

\begin{frame}[fragile=singleslide]{Use familiar sections, frames, and lists}
  \begin{columns}[T,onlytextwidth]
    \begin{column}{.54\textwidth}
      \small
\begin{verbatim}
\section{Results}
\begin{frame}{The main finding}
  \begin{itemize}
    \item State the result.
    \item Explain the evidence.
      \begin{itemize}
        \item Add supporting detail.
      \end{itemize}
  \end{itemize}
  \alert{The key takeaway.}
\end{frame}
\end{verbatim}
    \end{column}
    \begin{column}{.42\textwidth}
      \begin{itemize}
        \item State the result.
        \item Explain the evidence.
          \begin{itemize}
            \item Add supporting detail.
          \end{itemize}
      \end{itemize}
      \medskip
      \alert{The key takeaway.}
      \bigskip

      Use \texttt{enumerate} for numbered lists. Sections populate the
      sidebar automatically.
    \end{column}
  \end{columns}
\end{frame}

\begin{frame}[fragile]{Add Amurmaple's extra features}
  \begin{columns}[T,onlytextwidth]
    \begin{column}{.53\textwidth}
      \textbf{Inside a frame}
      {\small
\begin{verbatim}
\framesection{A short heading}
\boxalert{A point to remember}
\end{verbatim}
      }
      \framesection{A short heading}
      \boxalert{A point to remember}
      \bigskip

      The next pages show \texttt{information}, \texttt{remark}, and
      the standard Beamer boxes.
    \end{column}
    \begin{column}{.43\textwidth}
      \textbf{Between frames}
      {\small
\begin{verbatim}
\section{Results}
\sepframe[title={Results}]

% ... your result slides ...
\thanksframe{Questions?}
\end{verbatim}
      }
      These commands create complete divider and closing slides.
      \medskip

      For the full Amurmaple reference, see the bundled
      \texttt{beamer-amurmaple-doc.pdf}.
    \end{column}
  \end{columns}
\end{frame}
